
\subsection{Examining Zero-Shot Vulnerability Repair with Large Language Models}
\label{sec:vulnerability-repair}

\subsubsection{Motivation}

Pearce et al.~\cite{pearce2023repair} investigate whether general-purpose large language models (LLMs) can repair known software vulnerabilities without task-specific fine-tuning. The study aims to determine whether LLMs can reduce manual patching effort while maintaining both security and functional correctness.

\subsubsection{Methodology}

The authors evaluated seven LLMs across synthetic, hand-crafted, and real-world scenarios. They varied prompt content and model parameters to generate candidate patches. Static analysis and functional tests were used for synthetic software examples, while runtime sanitizers and regression tests were used for 12 real-world cases from ExtractFix. Existing developer patches were used to localize real-world bugs~\cite[Secs.~IV--V]{pearce2023repair}.

\subsubsection{Findings}

The models collectively generated successful repairs across all synthetic and hand-crafted scenarios. However, only 2.2\% of compilable candidates for out-of-bounds writes and 29.6\% for SQL injection were both functional and safe. In the real-world evaluation, eight of twelve cases yielded patches that passed the specified security and functional tests. Manual inspection nevertheless identified apparently successful patches that introduced other bugs~\cite[Secs.~IV--V]{pearce2023repair}.

\subsubsection{Conclusions}

Zero-shot vulnerability repair shows potential but was not sufficiently reliable for practical automated repair under the evaluated conditions. Improved bug localization, context selection, and stronger validation are needed. LLMs can assist developers in generating candidate patches, but their outputs require independent verification~\cite[Secs.~VI--IX]{pearce2023repair}.
